\documentclass[11pt,a4paper]{article}

\usepackage[margin=25mm]{geometry}
\usepackage[T1]{fontenc}
\usepackage[utf8]{inputenc}
\usepackage{amsmath,amssymb}
\usepackage{booktabs}
\usepackage{enumitem}
\usepackage{xcolor}
\usepackage{hyperref}
\usepackage{listings}
\usepackage{fancyhdr}
\usepackage{microtype}
\usepackage{parskip}

\definecolor{accent}{HTML}{244A73}
\definecolor{codebg}{HTML}{F4F6F8}
\definecolor{codegreen}{HTML}{276749}
\definecolor{codegray}{HTML}{5F6B76}

\hypersetup{
  colorlinks=true,
  linkcolor=accent,
  urlcolor=accent,
  pdftitle={PiPPINN Encoder-Decoder Design Notes},
  pdfsubject={Interpretability, reliability, PU learning, monotone maps, and KL stabilization}
}

\lstset{
  basicstyle=\ttfamily\small,
  backgroundcolor=\color{codebg},
  frame=single,
  rulecolor=\color{codebg},
  keywordstyle=\color{accent}\bfseries,
  commentstyle=\color{codegreen},
  stringstyle=\color{codegray},
  breaklines=true,
  showstringspaces=false,
  columns=fullflexible,
  xleftmargin=0.5em,
  xrightmargin=0.5em,
  aboveskip=0.8em,
  belowskip=0.8em
}

\pagestyle{fancy}
\fancyhf{}
\lhead{PiPPINN Model Design Notes}
\rhead{Encoder-Decoder Architecture}
\cfoot{\thepage}
\setlength{\headheight}{14pt}

\setlist[itemize]{leftmargin=1.5em,itemsep=0.25em,topsep=0.35em}
\setlist[enumerate]{leftmargin=1.7em,itemsep=0.35em,topsep=0.35em}

\title{\textcolor{accent}{PiPPINN Encoder-Decoder Design Notes}\\
\large Interpretable Evidence, Coverage-Aware PU Learning, Monotone Maps, and KL Stabilization}
\author{Technical discussion exported for offline reading}
\date{August 2026}

\begin{document}

\maketitle
\thispagestyle{fancy}

\begin{abstract}
This document explains several conceptual improvements to the PiPPINN
encoder-decoder while treating the upstream ESM-derived node features as fixed.
The focus is on the interpretation of additive decoder branches, evidence
reliability, positive-unlabeled learning under species-dependent coverage,
stable monotone mappings, and stable variational training. The unfinished
\texttt{HPO\_stage2.py} and \texttt{Evaluate\_Trained.py} workflows are outside
the scope of this discussion.
\end{abstract}

\tableofcontents
\newpage

\section{What ``not statistically identifiable or reliability-aware'' means}

These are two separate limitations. They do not imply that the biological
motivation of the transitivity or congruence branches is invalid.

\subsection{Statistical identifiability of branch contributions}

The current existence logit has the additive form

\begin{equation}
  \ell_{uv}
  = C_{\mathrm{trans}}(u,v)
  + C_{\mathrm{cong}}(u,v)
  + C_{\mathrm{MLP}}(u,v).
\end{equation}

Only the total is supervised. For example, a model could produce

\begin{center}
\begin{tabular}{lrr}
\toprule
Branch & Decomposition A & Decomposition B \\
\midrule
Transitivity & $+3$ & $+5$ \\
Congruence & $+2$ & $+1$ \\
Direct MLP & $-1$ & $-2$ \\
\midrule
Total logit & $+4$ & $+4$ \\
\bottomrule
\end{tabular}
\end{center}

The binary cross-entropy loss cannot distinguish these decompositions. This
ambiguity is possible because:

\begin{itemize}
  \item each branch can contain a trainable bias;
  \item the unrestricted MLP can learn signals correlated with homology or
        neighborhood evidence;
  \item only the final sum receives an interaction target.
\end{itemize}

The reported values are therefore exact computational contributions under the
current parameterization, but their allocation among branches may change across
training runs.

\subsubsection*{Minimal improvement: anchor each named branch}

Use one global baseline and make each named evidence function zero at a defined
neutral input:

\begin{equation}
  C_t(s) = f_t(s) - f_t(s_{t,\mathrm{neutral}}).
\end{equation}

The combined model becomes

\begin{equation}
  \ell_{uv}
  = b + C_{\mathrm{MLP}}(u,v)
  + g_t C_t(u,v) + g_c C_c(u,v),
\end{equation}

where $b$ is the only global prevalence baseline. For the deliberately
unnormalized transitivity score, a reasonable neutral reference is the expected
score of one random cosine comparison, approximately zero. This preserves the
intended interpretation that a larger number of supporting comparisons can add
evidence.

Anchoring removes arbitrary branch-baseline offsets. It does not remove every
possible redundancy with the MLP. Branch dropout, ablation tests, and
contribution-stability measurements across random seeds can quantify the
remaining redundancy.

\section{Evidence-reliability gates}

A reliability gate answers the question:

\begin{quote}
Given that an evidence branch produced a particular contribution, how strongly
should that contribution be trusted in this neighborhood context?
\end{quote}

It is not another interaction predictor. It scales evidence using information
about the support and precision of the evidence calculation.

Conceptually:

\begin{lstlisting}[language=Python]
transitivity_contribution = (
    transitivity_available
    * transitivity_reliability
    * centered_transitivity_evidence
)
\end{lstlisting}

The reliability is constrained to $[0,1]$. Multiplication attenuates both
positive and negative evidence when the underlying estimate is unreliable.

\subsection{Why a gate can help}

Similar aggregate scores can arise from contexts with different certainty:

\begin{itemize}
  \item a very small neighborhood with one strong match;
  \item many members of the smaller neighborhood consistently matching the
        larger neighborhood;
  \item a calculation based on a heavily truncated neighborhood;
  \item a score that changes substantially under a different sampled message
        context.
\end{itemize}

The unnormalized transitivity score represents the amount of supporting
similarity. The gate represents uncertainty in that amount. These are distinct
quantities.

\subsection{A simple analytical gate}

Let

\begin{equation}
  m = \min\bigl(|N(u)|,|N(v)|\bigr).
\end{equation}

A saturating count-reliability function is

\begin{equation}
  g_{\mathrm{count}} = \frac{m}{m+k},
\end{equation}

where $k$ is the amount of support that yields half reliability. If $k=5$,
then $m=1$, $m=5$, and $m=20$ yield approximately $0.17$, $0.50$, and $0.80$.

If the full known degree is available before neighborhood restriction, define

\begin{equation}
  q_u =
  \frac{|N_{\mathrm{retained}}(u)|}
       {\min\bigl(|N_{\mathrm{known}}(u)|,\mathrm{max\_neighbors}\bigr)},
\end{equation}

and analogously for $q_v$. A simple transitivity reliability is then

\begin{equation}
  g_t =
  \mathbf{1}_{\mathrm{available}}
  \frac{m}{m+k}\sqrt{q_uq_v}.
\end{equation}

An analytical gate is a good initial implementation. If it is learned, its
inputs should be limited to reliability variables such as counts, truncation
fractions, and sampling stability. Feeding latent similarity into the gate
would allow it to become an additional hidden interaction predictor.

An especially useful reliability measurement is context stability. A small
subset of candidate edges can occasionally be evaluated under two independently
sampled neighborhoods. Their contribution variance provides a direct estimate
of sampling uncertainty without doubling the cost for every edge.

\section{Complementary interactions and the direct decoder}

The deep direct decoder MLP does not assume latent similarity. Its sum and
elementwise-product representation is symmetric under endpoint exchange, while
the resulting function can remain arbitrary and non-monotonic. It can learn:

\begin{itemize}
  \item similar-role interactions;
  \item complementary-role interactions;
  \item incompatibilities;
  \item nonlinear combinations of latent properties.
\end{itemize}

A separate complementary expert is therefore unnecessary for predictive
expressiveness. Such an expert would only be useful if complementarity itself
needed to become another explicitly named explanatory mechanism.

There is one interpretive subtlety. The similarity branches and the MLP operate
on the same learned node latent. Gradients from transitivity and congruence pass
through the node encoder. Consequently, cosine similarity in that space is not
guaranteed to remain a pure measure of biological homology; it is a supervised
transformation of the ESM features. Unless this is validated against known
families or homology annotations, ``learned homology-like similarity'' is the
more precise description.

\section{Positive-unlabeled learning with coverage assumptions}

The intended reasoning is coherent:

\begin{itemize}
  \item an absent edge between well-covered proteins is stronger evidence of
        non-interaction;
  \item absence involving poorly covered proteins is weak evidence;
  \item node centrality is a proxy for local experimental coverage;
  \item species-level dataset quality changes how centrality should be
        interpreted.
\end{itemize}

Without negatives or external coverage data, true interaction probability and
observation probability cannot be learned independently. Some coverage
assumption is unavoidable. The improvement is to express that assumption as an
observation process rather than assigning an unlabeled pair an artificial
positive target.

\subsection{Separate biological truth from observation}

Let

\begin{itemize}
  \item $Y_{uv}=1$ mean that the true biological interaction exists;
  \item $O_{uv}=1$ mean that the interaction occurs in the observed dataset;
  \item $p_{uv}=P(Y_{uv}=1)$;
  \item $\pi_{uv}=P(O_{uv}=1\mid Y_{uv}=1)$ be the observation propensity.
\end{itemize}

Assuming recorded edges are genuine,

\begin{equation}
  P(O_{uv}=1) = \pi_{uv}p_{uv}.
\end{equation}

The negative log likelihood is

\begin{equation}
L_{uv}=
\begin{cases}
  -\log(\pi_{uv}p_{uv}), & O_{uv}=1,\\[4pt]
  -\log(1-\pi_{uv}p_{uv}), & O_{uv}=0.
\end{cases}
\end{equation}

If $\pi$ is fixed by the coverage model, $-\log\pi$ is constant for observed
positives. The effective interaction-model losses are

\begin{align}
  L_{\mathrm{positive}} &= -\log p,\\
  L_{\mathrm{unlabeled}} &= -\log(1-\pi p).
\end{align}

This has the desired behavior:

\begin{itemize}
  \item high $\pi$: a missing edge pushes $p$ downward strongly;
  \item low $\pi$: absence supplies little negative evidence;
  \item poor coverage does not imply that a missing pair is positive; it means
        that the observation is uninformative.
\end{itemize}

That last distinction matters. A soft target near $0.5$ tells the model that an
unobserved pair has high positive probability. The observation likelihood only
reduces the negative information supplied by its absence.

\subsection{Local coverage within each species}

Normalize degree within species so that high centrality means high coverage
relative to that organism:

\begin{equation}
  c_i^{(s)} =
  \operatorname{sigmoid}\left(
    a\frac{\log(1+d_i)-\operatorname{median}_s}
           {\operatorname{IQR}_s+\epsilon}
  \right).
\end{equation}

A percentile rank is a simpler alternative. A smooth pair-level coverage score
is

\begin{equation}
  c_{uv}^{(s)} = \sqrt{c_u^{(s)}c_v^{(s)}}.
\end{equation}

The geometric mean says that absence is highly reliable only if both proteins
are well covered. A sum encodes a different assumption: one highly studied
endpoint may be sufficient to make absence informative.

\subsection{Species-level coverage}

Introduce a species-specific coverage ceiling $g_s$:

\begin{equation}
  \pi_{uv}=g_s c_{uv}^{(s)}.
\end{equation}

A highly central Arabidopsis pair may have $c_{uv}\simeq1$, but a lower $g_s$
still prevents it from being treated as equally covered as a corresponding
human pair.

Without external information, the absolute $g_s$ values remain unidentified.
Practical options are:

\begin{enumerate}
  \item construct a fixed proxy from median degree, edges per protein, or
        another dataset-coverage statistic;
  \item specify plausible low, medium, and high coverage scenarios;
  \item check whether the important predictions remain stable across those
        scenarios;
  \item if $g_s$ is learned, constrain it strongly with priors and an externally
        selected reference species.
\end{enumerate}

Giving every species a freely learned propensity is risky: the model could
explain missing edges by reducing $g_s$, making interaction truth and observation
coverage inseparable.

\subsection{Simpler weighted-loss approximation}

If changing the likelihood is initially too invasive, use hard-zero unlabeled
targets and weight their loss by observation propensity:

\begin{equation}
  L_{\mathrm{unlabeled}}
  = -\pi_{uv}\log(1-p_{uv}).
\end{equation}

This approximates the observation likelihood when $p$ is not too large and is
more faithful to the coverage assumption than moving the unlabeled target
toward $0.5$.

\section{Monotone-map stabilization}

There are three distinct stabilization concerns: initialization, anchoring, and
path amplification.

\subsection{Initialization through inverse softplus}

The effective positive weight is currently

\begin{lstlisting}[language=Python]
weight = F.softplus(raw_weight) + epsilon
\end{lstlisting}

However,

\begin{equation}
  \operatorname{softplus}(0)\simeq0.693.
\end{equation}

Because all weights have the same sign, contributions from the many network
paths accumulate instead of canceling as they can under ordinary Xavier-style
initialization. Width-eight layers can therefore produce unexpectedly large
initial outputs.

To initialize the effective positive weight at $w_0$, initialize its raw value
with inverse softplus:

\begin{equation}
  r_0=\log\left(e^{w_0-\epsilon}-1\right).
\end{equation}

For example, an effective weight of approximately $0.1$ requires a raw value
near $-2.25$, not zero.

\subsection{Anchoring the output}

Softplus hidden activations are positive even for zero input. Small weights alone
therefore do not ensure zero output at neutral evidence. Explicit centering
provides that guarantee:

\begin{lstlisting}[language=Python]
contribution = monotone_map(score) - monotone_map(neutral_score)
\end{lstlisting}

This preserves monotonicity while preventing arbitrary branch-baseline drift.

\subsection{Controlling path amplification}

Because every path has the same sign, wide monotone networks can acquire large
slopes. Possible controls are:

\begin{itemize}
  \item normalize each positive weight row and learn a separate positive scale;
  \item penalize excessive derivatives;
  \item reduce the number of hidden units;
  \item replace the scalar monotone MLP with a one-dimensional monotonic
        piecewise-linear spline.
\end{itemize}

All relevant monotone inputs and outputs are scalar, so a monotonic spline is a
particularly transparent alternative: its knots directly show how an evidence
score changes the log odds.

The existence maps should be interpreted as logit contributions, not as
probabilities. They are unbounded and become probabilities only after the final
sigmoid.

\section{KL-divergence stabilization}

The current Gaussian parameterization can become unstable when

\begin{lstlisting}[language=Python]
std = torch.exp(0.5 * logvar)
\end{lstlisting}

is evaluated at extreme log variance. Large positive values generate large
samples and KL penalties; very negative values make the posterior nearly
deterministic.

\subsection{Stabilize the posterior scale}

One option is bounded log variance:

\begin{lstlisting}[language=Python]
logvar = logvar.clamp(-10, 6)
std = torch.exp(0.5 * logvar)
\end{lstlisting}

A smoother alternative predicts scale directly:

\begin{lstlisting}[language=Python]
std = F.softplus(raw_scale) + 1e-4
logvar = 2 * torch.log(std)
\end{lstlisting}

\subsection{Compute KL directly from log variance}

For a diagonal Gaussian posterior and standard normal prior:

\begin{lstlisting}[language=Python]
kl = 0.5 * (
    mu.square() + logvar.exp() - 1.0 - logvar
)
kl = kl.mean()
\end{lstlisting}

This avoids reconstructing log variance through
\texttt{log(std.pow(2) + epsilon)}.

\subsection{Warm up the KL coefficient}

Early in training, the decoder is poor. Applying the full KL penalty immediately
can encourage $\mu\rightarrow0$ and $\sigma\rightarrow1$ before the latent
variables become useful. Increase the coefficient gradually:

\begin{equation}
  \beta_t = \beta_{\mathrm{target}}
  \min\left(1,\frac{t}{T_{\mathrm{warmup}}}\right).
\end{equation}

A warm-up covering approximately the first 10--30 percent of training is a
reasonable starting point. Hyperparameter optimization can still choose the
final coefficient; the schedule determines how it is reached.

\subsection{Use free bits to preserve active dimensions}

Compute KL per latent dimension and do not penalize the first small information
allowance:

\begin{lstlisting}[language=Python]
kl_per_dim = 0.5 * (
    mu.square() + logvar.exp() - 1.0 - logvar
).mean(dim=0)

kl = torch.relu(kl_per_dim - free_bits).mean()
\end{lstlisting}

This lets each dimension carry a modest amount of information before the KL
regularizer pushes against it.

\subsection{Use deterministic routine validation}

For ordinary validation, use the posterior mean:

\begin{lstlisting}[language=Python]
nodes_latent = node_mu
\end{lstlisting}

For uncertainty estimation, draw several latent samples and report their mean
prediction and dispersion. A single new sample during each validation pass adds
noise without producing a reliable uncertainty estimate.

\subsection{Diagnostics}

Track the following during training:

\begin{itemize}
  \item mean KL per latent dimension;
  \item number of active latent dimensions;
  \item mean and range of posterior standard deviation;
  \item reconstruction or prediction loss;
  \item predictions obtained from $\mu$ versus sampled $z$.
\end{itemize}

If KL approaches zero while $\mu\simeq0$ and $\sigma\simeq1$, the posterior has
collapsed. Rapidly increasing KL or posterior scale indicates variance
instability.

\section{Recommended implementation order}

The most immediately useful sequence is:

\begin{enumerate}
  \item initialize monotone weights through inverse softplus;
  \item anchor named monotone contributions at explicit neutral inputs;
  \item introduce an analytical evidence-reliability gate;
  \item replace degree-based soft targets with a fixed or strongly constrained
        observation propensity;
  \item add KL warm-up and stable variance parameterization;
  \item use posterior means for routine validation and Monte Carlo samples only
        for uncertainty estimation;
  \item test contribution stability through branch ablations and repeated
        training seeds.
\end{enumerate}

These changes preserve the model's distinctive decoder logic. They primarily
improve the meaning, calibration, and stability of the existing mechanisms
rather than replacing them with a less interpretable architecture.

\end{document}
